\chapter{Interface commune}

\section{Organisation de l'écran}

L'espace authentifié est construit autour de quatre zones :

\begin{enumerate}[leftmargin=1.6em]
  \item la barre latérale, dont les entrées sont propres au rôle ;
  \item la barre supérieure avec recherche, notifications et compte ;
  \item le bandeau de page qui rappelle le contexte et la période ;
  \item l'espace de contenu avec vues, filtres, actions et détails.
\end{enumerate}

La barre latérale peut être réduite. Un badge rouge sur une entrée signale le
nombre d'éléments non lus ou nécessitant une attention dans ce module.

\section{Recherche globale}

La recherche supérieure permet de retrouver les objets autorisés pour le rôle :
stations et sessions pour un client, auxquels s'ajoutent notamment alertes et
travaux assignés pour les employés. Un résultat n'accorde jamais un droit
supplémentaire : l'ouverture reste soumise aux permissions et à l'organisation.

\section{Notifications}

La cloche regroupe les notifications récentes. Lorsqu'une notification concerne
une page précise, son entrée reste mise en évidence jusqu'à l'ouverture de
l'élément associé. Les badges des modules aident à distinguer la charge globale
des éléments propres à une page.

\begin{procedurebox}{Traiter une notification}
  \begin{enumerate}[leftmargin=1.6em]
    \item ouvrir la cloche ;
    \item lire le type, le niveau et la date ;
    \item ouvrir l'objet lié ;
    \item effectuer l'action métier requise ;
    \item vérifier que l'élément n'est plus marqué comme non lu.
  \end{enumerate}
\end{procedurebox}

\section{Vues, filtres et sélection}

Selon la page, un même ensemble de données peut être présenté en cartes, liste,
tableau, calendrier ou carte géographique. Les filtres limitent les résultats
sans modifier les données. Une sélection multiple n'est proposée que lorsqu'une
action globale existe réellement.

\begin{longtable}{|L{3.1cm}|L{5.1cm}|L{7.2cm}|}
  \hline
  \rowcolor{ctmint}\textbf{Contrôle} & \textbf{Usage} & \textbf{Bon réflexe} \\ \hline
  \endfirsthead
  \hline
  \rowcolor{ctmint}\textbf{Contrôle} & \textbf{Usage} & \textbf{Bon réflexe} \\ \hline
  \endhead
  Recherche locale & Filtrer la page courante & Utiliser une référence, un nom
  de station ou un terme significatif. \\ \hline
  Filtres & Restreindre par statut, date, rôle ou priorité & Réinitialiser les
  filtres avant de conclure qu'un élément est absent. \\ \hline
  Sélecteur de vue & Changer la représentation & Utiliser la carte pour la
  localisation, le tableau pour comparer et le calendrier pour planifier. \\ \hline
  Export & Produire un fichier CSV, JSON ou PDF & Choisir PDF pour un document
  lisible et CSV/JSON pour un traitement de données. \\ \hline
\end{longtable}

\section{Profil, paramètres et aide}

Le menu de l'avatar donne accès à :

\begin{itemize}[leftmargin=1.6em]
  \item \textbf{Profile} : informations personnelles, professionnelles et
  coordonnées ;
  \item \textbf{Settings} : préférences d'affichage, fuseau horaire, rayon de
  recherche et protection de l'accès ;
  \item \textbf{Help} : procédures adaptées au rôle et contact du support ;
  \item \textbf{Logout} : fermeture de la session.
\end{itemize}

Le fuseau horaire est une préférence du compte. Les dates affichées privilégient
ce fuseau local, alors que les événements techniques restent conservés de façon
cohérente côté serveur.

\section{Comprendre les principaux états}

\attention{L'état affiché d'une station n'est pas un simple champ saisi
manuellement. Il est recalculé à partir des connexions, \emph{heartbeats},
statuts des connecteurs, sessions, erreurs et maintenances.}

\begingroup
\small
\renewcommand{\arraystretch}{1.08}
\begin{tabularx}{\textwidth}{|L{3.2cm}|L{4.2cm}|Y|}
  \hline
  \rowcolor{ctmint}\textbf{Objet} & \textbf{État} & \textbf{Interprétation} \\ \hline
  Station & Disponible & Connecteur exploitable et communication valide. \\ \hline
  Station & En charge & Session active projetée depuis OCPP. \\ \hline
  Station & Non disponible & Réponse trop ancienne, défaut ou aucun connecteur
  exploitable. \\ \hline
  Station & Maintenance & Indisponibilité planifiée ou intervention en cours.
  \\ \hline
  Alerte & Nouvelle / En cours / Résolue & L'incident vient d'être détecté, est
  pris en charge ou a reçu une résolution. \\ \hline
  Intervention & Assignée / En cours / Terminée & Le technicien a reçu la tâche,
  l'exécute ou a remis son résultat. \\ \hline
\end{tabularx}
\endgroup
